% !TEX program = xelatex
\documentclass[a4paper,11pt,fontset=none]{ctexart}

\newcommand{\reporttitle}{A股低位主升浪全景：资金+技术+板块+标的}
\newcommand{\reportsubtitle}{2026年7月7日 | 三角形变盘 | 五大低位板块 | 逐日资金明细}
\newcommand{\reportkicker}{ASTOCK MARKET INTELLIGENCE}
\newcommand{\reportscope}{CHINA A-SHARES | FUND FLOW \& TECHNICALS}
\newcommand{\reportdate}{2026-07-07}
\newcommand{\reportdatacutoff}{2026-07-07 full-day closing data confirmed}
\newcommand{\reporttype}{Fund flow \& technical analysis}
\newcommand{\reportauthor}{AStock Agent Team}
\newcommand{\reporthouseview}{%
7月7日A股沪指收盘\textbf{3990.24点}（-1.26\%），正式失守4000点关口，为6月12日以来首次收于4000下方。全天成交\textbf{2.58万亿}（较前日骤缩5100亿），近4800股下跌，仅693只上涨（12.56\%）。技术面：日线三角形向下破位确认，MACD死叉扩散。资金面：全天主力净流出511--668亿（连续5日净流出），但北向资金分歧——一说全天净买入66亿（其中32.7亿进半导体设备），一说净卖出42.6亿（减持创新药/有色/白马）。核心确认：华天科技全天涨停+主力净流入\textbf{37.91亿}（特大单+37.18亿），中微公司大宗交易机构净买入3.89亿。半导体设备/先进封装为全天唯一正向板块，科创50唯一收红。结构性行情极致化，4000点已破，3970为下一关键支撑。%
}
\newcommand{\reportquality}{Full-day closing data (07/07 15:00); northbound/fund flow confirmed; broker reports cross-validated. High confidence.}
\newcommand{\reportdisclaimer}{本报告不构成任何投资建议。}
% Reusable IB-style preamble for XeLaTeX equity research reports
% Usage: % Reusable IB-style preamble for XeLaTeX equity research reports
% Usage: % Reusable IB-style preamble for XeLaTeX equity research reports
% Usage: \input{preamble} after \documentclass[a4paper,11pt,openany,fontset=none]{ctexrep}

% ============ 页面布局 ============
\usepackage[top=2cm, bottom=2cm, left=2cm, right=2cm]{geometry}
\setlength{\parskip}{0.3em}
\setlength{\parindent}{2em}
\renewcommand{\baselinestretch}{1.15}

% ============ 字体（macOS） ============
\setCJKmainfont[BoldFont={Heiti SC}]{STSong}
\setCJKsansfont{Heiti SC}
\setCJKmonofont{Heiti SC}
\setmainfont{Times New Roman}

% ============ 颜色（IB Navy/Grey palette） ============
\usepackage{xcolor}
\definecolor{navy}{HTML}{003366}
\definecolor{steelgrey}{HTML}{4A5568}
\definecolor{lightgrey}{HTML}{F7FAFC}
\definecolor{riskred}{HTML}{C53030}
\definecolor{riskamber}{HTML}{D69E2E}
\definecolor{riskgreen}{HTML}{38A169}
\definecolor{nvgreen}{HTML}{76B900}

% ============ 表格 ============
\usepackage{booktabs}
\usepackage{longtable}
\usepackage{tabularx}
\usepackage{multirow}
\usepackage{array}
\newcolumntype{Y}{>{\centering\arraybackslash}X}
\newcolumntype{L}[1]{>{\raggedright\arraybackslash}p{#1}}
\newcolumntype{C}[1]{>{\centering\arraybackslash}p{#1}}
\newcolumntype{R}[1]{>{\raggedleft\arraybackslash}p{#1}}

% ============ 图形 ============
\usepackage{tikz}
\usetikzlibrary{shapes,arrows.meta,positioning,calc,backgrounds,fit}
\usepackage{pgfplots}
\usepgfplotslibrary{polar}
\pgfplotsset{compat=1.18}
\usepackage[edges]{forest}

% ============ 页眉页脚（报告标题和日期需在main.tex中覆盖） ============
\usepackage{fancyhdr}
\pagestyle{fancy}
\fancyhf{}
\fancyhead[L]{\small\color{steelgrey} \reporttitle}
\fancyhead[R]{\small\color{steelgrey} \reportdate}
\fancyfoot[C]{\small\color{steelgrey} \thepage}
\fancyfoot[R]{\tiny\color{steelgrey} 本报告不构成任何投资建议}
\renewcommand{\headrulewidth}{0.4pt}
\renewcommand{\footrulewidth}{0.2pt}

% ============ 章节格式 ============
\usepackage{titlesec}
\titleformat{\chapter}[display]
  {\normalfont\huge\bfseries\color{navy}}
  {\chaptertitlename\ \thechapter}{20pt}{\Huge}
\titleformat{\section}
  {\normalfont\Large\bfseries\color{navy}}
  {\thesection}{1em}{}
\titleformat{\subsection}
  {\normalfont\large\bfseries\color{steelgrey}}
  {\thesubsection}{1em}{}

% ============ 超链接 ============
\usepackage{hyperref}
\hypersetup{
  colorlinks=true,
  linkcolor=navy,
  urlcolor=navy,
  citecolor=navy,
}

% ============ 其他工具 ============
\usepackage{enumitem}
\setlist{nosep, leftmargin=2em}
\usepackage{fontawesome5}
\usepackage{amssymb}
\usepackage{pifont}
\usepackage{tcolorbox}
\tcbuselibrary{skins,breakable}
\usepackage{caption}
\captionsetup{font=small, skip=4pt}
\usepackage{float}
\setlength{\textfloatsep}{8pt plus 2pt minus 2pt}
\setlength{\floatsep}{6pt plus 2pt minus 2pt}
\setlength{\intextsep}{6pt plus 2pt minus 2pt}
\setlength{\abovecaptionskip}{4pt}
\setlength{\belowcaptionskip}{2pt}

% ============ 自定义命令 ============
\newcommand{\rmark}{\textcolor{riskred}{\ding{108}}}
\newcommand{\amark}{\textcolor{riskamber}{\ding{108}}}
\newcommand{\gmark}{\textcolor{riskgreen}{\ding{108}}}
\newcommand{\starfive}{$\bigstar\bigstar\bigstar\bigstar\bigstar$}
\newcommand{\starfour}{$\bigstar\bigstar\bigstar\bigstar$}
\newcommand{\starthree}{$\bigstar\bigstar\bigstar$}

% 信息框
\newtcolorbox{keyinsight}[1][]{
  colback=lightgrey, colframe=navy,
  fonttitle=\bfseries, title={#1}, breakable
}
\newtcolorbox{riskbox}[1][]{
  colback=red!5, colframe=riskred,
  fonttitle=\bfseries, title={#1}, breakable
}
 after \documentclass[a4paper,11pt,openany,fontset=none]{ctexrep}

% ============ 页面布局 ============
\usepackage[top=2cm, bottom=2cm, left=2cm, right=2cm]{geometry}
\setlength{\parskip}{0.3em}
\setlength{\parindent}{2em}
\renewcommand{\baselinestretch}{1.15}

% ============ 字体（macOS） ============
\setCJKmainfont[BoldFont={Heiti SC}]{STSong}
\setCJKsansfont{Heiti SC}
\setCJKmonofont{Heiti SC}
\setmainfont{Times New Roman}

% ============ 颜色（IB Navy/Grey palette） ============
\usepackage{xcolor}
\definecolor{navy}{HTML}{003366}
\definecolor{steelgrey}{HTML}{4A5568}
\definecolor{lightgrey}{HTML}{F7FAFC}
\definecolor{riskred}{HTML}{C53030}
\definecolor{riskamber}{HTML}{D69E2E}
\definecolor{riskgreen}{HTML}{38A169}
\definecolor{nvgreen}{HTML}{76B900}

% ============ 表格 ============
\usepackage{booktabs}
\usepackage{longtable}
\usepackage{tabularx}
\usepackage{multirow}
\usepackage{array}
\newcolumntype{Y}{>{\centering\arraybackslash}X}
\newcolumntype{L}[1]{>{\raggedright\arraybackslash}p{#1}}
\newcolumntype{C}[1]{>{\centering\arraybackslash}p{#1}}
\newcolumntype{R}[1]{>{\raggedleft\arraybackslash}p{#1}}

% ============ 图形 ============
\usepackage{tikz}
\usetikzlibrary{shapes,arrows.meta,positioning,calc,backgrounds,fit}
\usepackage{pgfplots}
\usepgfplotslibrary{polar}
\pgfplotsset{compat=1.18}
\usepackage[edges]{forest}

% ============ 页眉页脚（报告标题和日期需在main.tex中覆盖） ============
\usepackage{fancyhdr}
\pagestyle{fancy}
\fancyhf{}
\fancyhead[L]{\small\color{steelgrey} \reporttitle}
\fancyhead[R]{\small\color{steelgrey} \reportdate}
\fancyfoot[C]{\small\color{steelgrey} \thepage}
\fancyfoot[R]{\tiny\color{steelgrey} 本报告不构成任何投资建议}
\renewcommand{\headrulewidth}{0.4pt}
\renewcommand{\footrulewidth}{0.2pt}

% ============ 章节格式 ============
\usepackage{titlesec}
\titleformat{\chapter}[display]
  {\normalfont\huge\bfseries\color{navy}}
  {\chaptertitlename\ \thechapter}{20pt}{\Huge}
\titleformat{\section}
  {\normalfont\Large\bfseries\color{navy}}
  {\thesection}{1em}{}
\titleformat{\subsection}
  {\normalfont\large\bfseries\color{steelgrey}}
  {\thesubsection}{1em}{}

% ============ 超链接 ============
\usepackage{hyperref}
\hypersetup{
  colorlinks=true,
  linkcolor=navy,
  urlcolor=navy,
  citecolor=navy,
}

% ============ 其他工具 ============
\usepackage{enumitem}
\setlist{nosep, leftmargin=2em}
\usepackage{fontawesome5}
\usepackage{amssymb}
\usepackage{pifont}
\usepackage{tcolorbox}
\tcbuselibrary{skins,breakable}
\usepackage{caption}
\captionsetup{font=small, skip=4pt}
\usepackage{float}
\setlength{\textfloatsep}{8pt plus 2pt minus 2pt}
\setlength{\floatsep}{6pt plus 2pt minus 2pt}
\setlength{\intextsep}{6pt plus 2pt minus 2pt}
\setlength{\abovecaptionskip}{4pt}
\setlength{\belowcaptionskip}{2pt}

% ============ 自定义命令 ============
\newcommand{\rmark}{\textcolor{riskred}{\ding{108}}}
\newcommand{\amark}{\textcolor{riskamber}{\ding{108}}}
\newcommand{\gmark}{\textcolor{riskgreen}{\ding{108}}}
\newcommand{\starfive}{$\bigstar\bigstar\bigstar\bigstar\bigstar$}
\newcommand{\starfour}{$\bigstar\bigstar\bigstar\bigstar$}
\newcommand{\starthree}{$\bigstar\bigstar\bigstar$}

% 信息框
\newtcolorbox{keyinsight}[1][]{
  colback=lightgrey, colframe=navy,
  fonttitle=\bfseries, title={#1}, breakable
}
\newtcolorbox{riskbox}[1][]{
  colback=red!5, colframe=riskred,
  fonttitle=\bfseries, title={#1}, breakable
}
 after \documentclass[a4paper,11pt,openany,fontset=none]{ctexrep}

% ============ 页面布局 ============
\usepackage[top=2cm, bottom=2cm, left=2cm, right=2cm]{geometry}
\setlength{\parskip}{0.3em}
\setlength{\parindent}{2em}
\renewcommand{\baselinestretch}{1.15}

% ============ 字体（macOS） ============
\setCJKmainfont[BoldFont={Heiti SC}]{STSong}
\setCJKsansfont{Heiti SC}
\setCJKmonofont{Heiti SC}
\setmainfont{Times New Roman}

% ============ 颜色（IB Navy/Grey palette） ============
\usepackage{xcolor}
\definecolor{navy}{HTML}{003366}
\definecolor{steelgrey}{HTML}{4A5568}
\definecolor{lightgrey}{HTML}{F7FAFC}
\definecolor{riskred}{HTML}{C53030}
\definecolor{riskamber}{HTML}{D69E2E}
\definecolor{riskgreen}{HTML}{38A169}
\definecolor{nvgreen}{HTML}{76B900}

% ============ 表格 ============
\usepackage{booktabs}
\usepackage{longtable}
\usepackage{tabularx}
\usepackage{multirow}
\usepackage{array}
\newcolumntype{Y}{>{\centering\arraybackslash}X}
\newcolumntype{L}[1]{>{\raggedright\arraybackslash}p{#1}}
\newcolumntype{C}[1]{>{\centering\arraybackslash}p{#1}}
\newcolumntype{R}[1]{>{\raggedleft\arraybackslash}p{#1}}

% ============ 图形 ============
\usepackage{tikz}
\usetikzlibrary{shapes,arrows.meta,positioning,calc,backgrounds,fit}
\usepackage{pgfplots}
\usepgfplotslibrary{polar}
\pgfplotsset{compat=1.18}
\usepackage[edges]{forest}

% ============ 页眉页脚（报告标题和日期需在main.tex中覆盖） ============
\usepackage{fancyhdr}
\pagestyle{fancy}
\fancyhf{}
\fancyhead[L]{\small\color{steelgrey} \reporttitle}
\fancyhead[R]{\small\color{steelgrey} \reportdate}
\fancyfoot[C]{\small\color{steelgrey} \thepage}
\fancyfoot[R]{\tiny\color{steelgrey} 本报告不构成任何投资建议}
\renewcommand{\headrulewidth}{0.4pt}
\renewcommand{\footrulewidth}{0.2pt}

% ============ 章节格式 ============
\usepackage{titlesec}
\titleformat{\chapter}[display]
  {\normalfont\huge\bfseries\color{navy}}
  {\chaptertitlename\ \thechapter}{20pt}{\Huge}
\titleformat{\section}
  {\normalfont\Large\bfseries\color{navy}}
  {\thesection}{1em}{}
\titleformat{\subsection}
  {\normalfont\large\bfseries\color{steelgrey}}
  {\thesubsection}{1em}{}

% ============ 超链接 ============
\usepackage{hyperref}
\hypersetup{
  colorlinks=true,
  linkcolor=navy,
  urlcolor=navy,
  citecolor=navy,
}

% ============ 其他工具 ============
\usepackage{enumitem}
\setlist{nosep, leftmargin=2em}
\usepackage{fontawesome5}
\usepackage{amssymb}
\usepackage{pifont}
\usepackage{tcolorbox}
\tcbuselibrary{skins,breakable}
\usepackage{caption}
\captionsetup{font=small, skip=4pt}
\usepackage{float}
\setlength{\textfloatsep}{8pt plus 2pt minus 2pt}
\setlength{\floatsep}{6pt plus 2pt minus 2pt}
\setlength{\intextsep}{6pt plus 2pt minus 2pt}
\setlength{\abovecaptionskip}{4pt}
\setlength{\belowcaptionskip}{2pt}

% ============ 自定义命令 ============
\newcommand{\rmark}{\textcolor{riskred}{\ding{108}}}
\newcommand{\amark}{\textcolor{riskamber}{\ding{108}}}
\newcommand{\gmark}{\textcolor{riskgreen}{\ding{108}}}
\newcommand{\starfive}{$\bigstar\bigstar\bigstar\bigstar\bigstar$}
\newcommand{\starfour}{$\bigstar\bigstar\bigstar\bigstar$}
\newcommand{\starthree}{$\bigstar\bigstar\bigstar$}

% 信息框
\newtcolorbox{keyinsight}[1][]{
  colback=lightgrey, colframe=navy,
  fonttitle=\bfseries, title={#1}, breakable
}
\newtcolorbox{riskbox}[1][]{
  colback=red!5, colframe=riskred,
  fonttitle=\bfseries, title={#1}, breakable
}


\begin{document}

\begin{center}
{\sffamily\small\bfseries\color{steelgrey} \reportkicker}\\[3pt]
{\LARGE\bfseries\color{navy} \reporttitle}\\[5pt]
{\large\color{steelgrey} \reportsubtitle}\\[4pt]
{\small\color{mutedgrey} \reportscope \quad|\quad \reportdate \quad|\quad \reportdatacutoff}
\end{center}
\vspace{0.7em}

\begin{dashboardbox}[Decision Snapshot]
\begin{houseviewbox}[House View]
\reporthouseview
\end{houseviewbox}
\begin{sourcequalitybox}[Data Quality]
\reportquality
\end{sourcequalitybox}
\end{dashboardbox}

%% ================================================================
\section{大盘技术面：三角形末端破位}

\subsection*{指数表现（7月7日收盘）}

\begin{tabularx}{\textwidth}{L{3cm}L{2.5cm}L{2cm}X}
\toprule
\textbf{指数} & \textbf{收盘点位} & \textbf{涨跌幅} & \textbf{信号} \\
\midrule
上证指数 & \textbf{3990.24} & \textbf{-1.26\%} & 失守4000点（6/12以来首次） \\
深证成指 & 15225.11 & -1.24\% & 同步走弱 \\
创业板指 & 3911.91 & -0.94\% & 相对抗跌 \\
科创50 & --- & \textbf{+0.28\%} & 唯一收红（半导体设备撑） \\
北证50 & --- & -0.79\% & 弱 \\
\bottomrule
\end{tabularx}

\textbf{全天数据}：成交额\textbf{2.58万亿}（较前日3.09万亿骤缩5100亿），上涨693只（12.56\%），下跌4797只。缩量+普跌=恐慌衰竭但无增量。

\subsection*{核心技术形态}

\begin{exhibitbox}[Exhibit 1: 上证日线收敛三角形——变盘窗口已到]
\begin{tabularx}{\textwidth}{L{3.5cm}X}
\toprule
\textbf{技术要素} & \textbf{状态} \\
\midrule
形态 & \textbf{日线收敛三角形}（6月以来持续20+交易日），已运行至三角形末端 \\
小三角下沿 & 7/7上午已\textbf{破位}（盘中跌破4013点支撑） \\
大三角下沿 & \textbf{3970点}——自4258点顶部以来的关键趋势线 \\
均线 & 5/10/20/60日均线高度粘合缠绕，压力位4060--4120点 \\
MACD & 已死叉（7/3跌破预测指标线4069确认），DIF向下发散 \\
量能 & 连续4日缩量（3.09万亿$\to$半日1.63万亿），场外增量资金缺失 \\
震荡中枢 & 4040点一线，区间宽度仅60点 \\
7/6收盘形态 & 长下影锤头线（4000点获承接），当日验证支撑有效 \\
7/7午盘形态 & 向下突破小三角，盘中触及3999.03 \\
\bottomrule
\end{tabularx}
\end{exhibitbox}

\subsection*{关键位置图}

\begin{tabularx}{\textwidth}{L{3cm}L{2cm}X}
\toprule
\textbf{价位} & \textbf{性质} & \textbf{含义} \\
\midrule
4120点 & 三角形上轨+均线压力 & 反弹上限，不放量突破则继续震荡 \\
4060点 & 短期压力位 & 多条均线汇聚 \\
4040点 & 震荡中枢 & 多空平衡点 \\
4028点 & 短期支撑 & 跌破进入下一台阶 \\
4013点 & 短期强支撑 & 周一已验证承接，7/7已破 \\
\textbf{4000点} & \textbf{心理关口} & 政策/机构维稳线 \\
\textbf{3970点} & \textbf{大三角下沿} & 若失守$\to$自4258点顶部以来2浪C开启 \\
3794点 & 3月底部 & 极限防守位 \\
\bottomrule
\end{tabularx}

\subsection*{技术面结论}

\begin{riskbox}[变盘方向判断]
\begin{itemize}
  \item \textbf{基准情景（概率60\%）}：4000--3970区间获承接，继续震荡磨底1--2周，等待中报业绩催化后向上选择方向。
  \item \textbf{偏空情景（概率30\%）}：有效跌破3970$\to$开启中期调整C浪，目标3800--3850。
  \item \textbf{偏多情景（概率10\%）}：增量资金突然入场+放量站稳4060$\to$三角形向上突破。
\end{itemize}
\end{riskbox}

%% ================================================================
\section{资金面全景：六维度拆解}

\subsection*{维度一：主力资金（7/7全天收盘）}

\begin{exhibitbox}[Exhibit 2: 7/7全天主力资金——连续5日净流出]
\begin{tabularx}{\textwidth}{L{3cm}X}
\toprule
\textbf{方向} & \textbf{数据} \\
\midrule
全天主力净流出 & \textbf{-511亿}（证券时报）/ -668亿（数据宝），\textbf{连续5个交易日净流出} \\
特大单净流出 & -412.75亿元（1506股流入 vs 2939股流出） \\
流入行业（仅8个） & 游戏+12.6亿、银行+6.47亿、光伏设备+4.73亿 \\
流出行业 & 电子、通信、医药为主 \\
\bottomrule
\end{tabularx}
\end{exhibitbox}

\textbf{7/7全天主力净流入TOP个股}（收盘确认）：
\begin{enumerate}
  \item \textbf{华天科技}：主力净流入\textbf{+37.91亿}（特大单+37.18亿），涨停，成交120亿，龙虎榜净买入15.77亿
  \item 东山精密：净流入+13.3亿
  \item 中际旭创：净流入+11.1亿（回调抄底）
  \item 海光信息：净流入+8.44亿
  \item 有研新材：涨停
  \item 长电科技：封装权重
  \item 中微公司：+3.07亿，大宗交易机构净买入3.89亿
\end{enumerate}

\textbf{关键变化}（vs午盘）：全天主力从午盘``结构性净流入124亿''恶化为全天``净流出511亿''——\textbf{下午资金大幅流出}。但华天科技从半日28.85亿扩大到全天37.91亿（封板锁定），半导体设备方向逆势确认。

\subsection*{维度二：北向资金（收盘确认——数据存在口径分歧）}

\begin{tabularx}{\textwidth}{L{3cm}X}
\toprule
\textbf{数据源} & \textbf{7/7全天数据} \\
\midrule
口径A（收评/板块） & 全天净买入\textbf{66亿}，其中32.7亿全部进入半导体设备/光刻/硅片/封装 \\
口径B（微博/个别） & 全天净卖出42.6亿（减持创新药/有色/白马，小幅加仓半导体上游） \\
深股通成交TOP3 & 中际旭创57.26亿、新易盛42.40亿、\textbf{北方华创30.62亿} \\
沪股通成交TOP3 & 兆易创新36.14亿、寒武纪28.90亿、\textbf{中微公司27.08亿} \\
\bottomrule
\end{tabularx}

\textbf{核心确认}：无论整体口径如何，\textbf{半导体设备/材料是北向唯一确认加仓方向}（北方华创30.62亿成交+中微27.08亿成交位列前排）。创新药/有色/消费被减持。

\subsection*{维度三：两融杠杆资金}

\begin{tabularx}{\textwidth}{L{3.5cm}X}
\toprule
\textbf{指标} & \textbf{数据} \\
\midrule
两融余额（7/6） & 29,926.14亿元（占流通市值2.87\%） \\
单日变动 & \textbf{-107.23亿}（减少，杠杆资金降温） \\
融资余额变动 & -104.96亿（融资减少$\to$多头信心下降） \\
科创板两融 & 4104.49亿，\textbf{连续3日减少}（-40.21亿） \\
ETF两融 & 1163.29亿，+10.34亿（ETF端小幅净加杠杆） \\
\bottomrule
\end{tabularx}

\textbf{判断}：杠杆资金净减仓107亿/日，但无平仓踩踏迹象。科创板连续3日去杠杆（高位科技兑现），ETF端微幅加杠杆（机构借ETF布局低位板块）。整体杠杆水平平稳，不构成系统性风险。

\subsection*{维度四：解禁与减持压力}

\begin{tabularx}{\textwidth}{L{3.5cm}X}
\toprule
\textbf{指标} & \textbf{数据} \\
\midrule
7月全月解禁 & \textbf{2499亿元}（146家公司，169亿股） \\
7/6--7/10单周 & 超1000亿（其中屹唐股份单只653.62亿） \\
集中赛道 & 算力、半导体、通信（上半年翻倍的热门板块） \\
解禁/流通比 & 多只龙头>30\%，产业资本套现意愿极强 \\
同比 & -48\%（比去年同期少，但结构集中） \\
\bottomrule
\end{tabularx}

\textbf{判断}：解禁是高位科技赛道资金出逃的核心催化。产业股东持仓成本极低+上半年股价翻倍$\to$锁定套现。这解释了为什么资金从高位系统性迁出。\textbf{低位板块（设备/创新药/军工）解禁压力极小}，是资金避风港。

\subsection*{维度五：ETF与公募申赎}

\begin{tabularx}{\textwidth}{L{3.5cm}X}
\toprule
\textbf{数据} & \textbf{判断} \\
\midrule
公募申赎 & 正常，无大规模赎回 \\
军工ETF & 连续3日净申购1.09亿+ \\
半导体设备ETF & 持续净申购 \\
存储ETF & 大幅净申购（业绩催化） \\
高位算力ETF & 部分净赎回 \\
\bottomrule
\end{tabularx}

\subsection*{维度六：量能结构}

\begin{tabularx}{\textwidth}{L{3cm}X}
\toprule
\textbf{日期} & \textbf{全天成交} \\
\midrule
7/3（周四） & 3.99万亿（恐慌放量杀跌） \\
7/4（周五） & 3.18万亿（缩量） \\
7/6（周一） & 3.09万亿（继续缩量） \\
\textbf{7/7（周二）} & \textbf{2.58万亿}（骤缩5100亿，连续4日缩量） \\
\bottomrule
\end{tabularx}

\textbf{量能判断（收盘更新）}：
\begin{itemize}
  \item 单日缩量5100亿是极端信号——场外资金彻底观望，场内仅存量搬砖。
  \item 从3.99万亿$\to$2.58万亿，4个交易日成交骤降35\%，已接近``地量''水平。
  \item \textbf{缩量至极致本身是见底条件之一}：卖无可卖$\to$抛压枯竭$\to$等待催化反转。
  \item 但反转需要量能重新放大至3万亿+配合，否则仅是``温水煮蛙''阴跌。
\end{itemize}

%% ================================================================
\section{板块资金流向热力图}

\begin{exhibitbox}[Exhibit 3: 7月初资金流向全景（按确定性排序）]
\begin{tabularx}{\textwidth}{L{0.5cm}L{3cm}L{2.5cm}L{2.5cm}X}
\toprule
 & \textbf{板块} & \textbf{7月资金方向} & \textbf{北向态度} & \textbf{阶段判断} \\
\midrule
1 & 先进封装 & 主力大幅流入 & 加仓 & \textbf{当日最强}（华天28亿） \\
2 & 半导体设备/材料 & 北向110亿+稳步吸纳 & 重仓加码 & 底部蓄势末期 \\
3 & 存储芯片 & 业绩催化流入 & 加仓 & 高位震荡+分化 \\
4 & 国产算力/信创 & 逆势抗跌有承接 & 加仓 & 低位结构性机会 \\
5 & 创新药 & 连续3日流入>10亿 & 中性 & 已走出主升（7/7有兑现） \\
6 & 军工 & ETF连续申购+社保 & 逆势加仓 & 底部建仓中 \\
7 & 氟化工/电子特气 & 稳步流入 & 加仓 & 涨价确认期 \\
\midrule
8 & AI算力/光模块 & \textbf{大幅流出290亿+} & 减持 & 高位兑现进行中 \\
9 & 题材小票 & 出逃 & 减持 & 回避 \\
10 & 面板/消费电子 & 分化 & 中性 & 等中报确认 \\
\bottomrule
\end{tabularx}
\end{exhibitbox}

%% ================================================================
\section{龙虎榜与机构动向（7/7确认）}

\begin{tabularx}{\textwidth}{L{3cm}L{2cm}X}
\toprule
\textbf{个股} & \textbf{金额} & \textbf{机构动作} \\
\midrule
华天科技 & +28.85亿 & 封测龙头，机构巨量抢筹 \\
长电科技 & +25.59亿 & 封装权重，机构加仓 \\
蔚蓝锂芯 & +15.29亿 & 机构龙虎榜大额买入 \\
金安国纪 & +6.87亿 & 3家机构集体抄底（大跌-9.5\%后） \\
惠科股份 & +5.68亿 & 机构买入 \\
雷赛智能 & +1.89亿 & 机器人伺服龙头，4连板 \\
\bottomrule
\end{tabularx}

\textbf{机构行为特征}：
\begin{itemize}
  \item 偏好抄底回调个股（金安国纪大跌后3机构集体加仓）。
  \item 高涨个股（紫光股份、中石科技）机构反而集中卖出——高位止盈明确。
  \item 北向同步抄底超跌（中钨高新、永鼎股份）+ 加仓算力（中国长城、菲菱科思）。
\end{itemize}

%% ================================================================
\section{市场微观结构分析}

\begin{tabularx}{\textwidth}{L{4cm}X}
\toprule
\textbf{微观指标} & \textbf{状态与含义（收盘确认）} \\
\midrule
涨跌家数 & 上涨693只（12.56\%） vs \textbf{下跌4797只} \\
成交额 & 2.58万亿（骤缩5100亿） \\
涨停 & 华天科技涨停，有研新材涨停，科创半导体设备ETF+3\% \\
跌停 & 多只（有色、题材小票） \\
科创50 & \textbf{唯一收红宽基}（+0.28\%）——半导体设备撑住 \\
市场情绪 & 极度悲观，但缩量至极致=卖无可卖 \\
\bottomrule
\end{tabularx}

%% ================================================================
\section{综合判断与操作策略}

\subsection*{大盘趋势判断}

\begin{tabularx}{\textwidth}{L{2cm}L{1.5cm}X}
\toprule
\textbf{情景} & \textbf{概率} & \textbf{触发条件与路径} \\
\midrule
震荡磨底 & 60\% & 3970--4060区间震荡1--2周，缩量企稳后等中报催化向上 \\
中期调整 & 30\% & 有效跌破3970+放量$\to$C浪下跌至3800--3850 \\
V形反转 & 10\% & 政策/外部突发利好+增量资金入场+放量站稳4060 \\
\bottomrule
\end{tabularx}

\subsection*{资金面核心结论}

\begin{enumerate}
  \item \textbf{高位出逃不可逆}：7月解禁2499亿+上半年利润兑现+机构半年考核=高位科技持续失血。至少延续至7月下旬。
  \item \textbf{低位承接确定}：北向98亿逆势流入+龙虎榜机构/社保扫货+ETF持续净申购=低位板块建仓确认。
  \item \textbf{增量资金缺失}：两融减仓+成交萎缩+场外观望=存量博弈，指数级行情暂时不具备条件。
  \item \textbf{结构性为王}：板块间资金流动远大于大盘方向性，跟对方向比判断指数重要10倍。
\end{enumerate}

\subsection*{操作策略}

\begin{tabularx}{\textwidth}{L{3cm}X}
\toprule
\textbf{类型} & \textbf{策略} \\
\midrule
仓位管理 & 整体仓位控制在5--7成，留3成现金应对破位风险 \\
方向选择 & \textbf{做多低位资金流入方向}：先进封装>半导体设备>军工>氟化工 \\
回避方向 & 上半年翻倍的高位科技/光模块/算力题材（解禁+兑现双杀） \\
右侧确认 & 若沪指放量站稳4060$\to$可加仓至8成 \\
止损线 & 沪指有效跌破3970（收盘确认）$\to$整体降至3成以下 \\
\bottomrule
\end{tabularx}

\subsection*{时间节奏}

\begin{itemize}
  \item \textbf{本周（7/7--7/11）}：三角形方向选择窗口，后半周大概率明确。重点观察：(1) 3970是否有效守住；(2) 券商板块能否异动（唯一能搅动做多气氛的权重）。
  \item \textbf{7月中旬}：中报密集预告期，业绩确定性强的标的将获增量资金。
  \item \textbf{7月下旬}：高低切换完成，资金可能重新回流经过充分调整的硬科技龙头。
\end{itemize}

%% ================================================================
\newpage
\section{五大低位蓄势板块：基本面+资金面+标的}

基于前述资金流向数据，以下五大板块已获得机构/北向/ETF多路资金确认性流入，处于低位蓄势或主升浪早期。每个板块附带核心驱动、资金量化证据、以及重点标的。

\subsection*{板块一：先进封装（当日最强，7/7半日吸金54亿+）}

\textbf{资金证据}：华天科技7/7半日+28.85亿（全市场第一）+长电科技+25.59亿。7/6华天主力净流入+19.94亿。封测板块单日合计吸金断层领先全市场。

\textbf{驱动逻辑}：AI服务器/存储芯片封装需求爆发+中报预告净利翻倍+大基金三期扶持Chiplet/2.5D/3D混合键合。

\textbf{重点标的}：
\begin{itemize}
  \item \textbf{华天科技}（002185）：7/7涨停，半日主力+28.85亿；今年6次龙虎榜；中报净利翻倍。
  \item \textbf{长电科技}（600584）：7/7主力+25.59亿，封装权重股。
  \item 通富微电（002156）、深科技（000021）。
\end{itemize}

\begin{tabularx}{\textwidth}{L{2.5cm}L{2cm}L{2cm}L{3cm}X}
\toprule
\textbf{标的} & \textbf{当前价} & \textbf{目标价} & \textbf{机构} & \textbf{空间} \\
\midrule
华天科技 & $\sim$21.6元 & 23.5元 & 花旗（6/23买入） & +9\% \\
华天科技 & & 25--28元 & 乐观估值 & +16--30\% \\
长电科技 & $\sim$97元 & 125元 & 高盛（7/2上调） & +29\% \\
长电科技 & & 110元 & 摩根大通（6/16买入） & +13\% \\
\bottomrule
\end{tabularx}

\subsection*{板块二：半导体设备/材料（北向110亿+锁仓，底部蓄势末期）}

\textbf{资金证据}：
\begin{itemize}
  \item 北向7月累计板块净流入\textbf{110亿+}，单只4--19亿/月。
  \item 北方华创：北向持仓9531万股/13.16\%流通+大基金一期4.80\%+二期+482家公募8.88\%。
  \item 中微公司：北向Q1\textbf{增持1069万股}+大基金3.49\%+设备ETF新进。
  \item 设备ETF（561980）10日累计吸金22亿+，7/7逆势涨2\%。
  \item 北方华创7/3跌停次日主力暗吸5.53亿；7/6游资15亿逆势进场。
\end{itemize}

\textbf{驱动逻辑}：国内晶圆厂扩产订单锁定至2027+国产化率20\%$\to$40\%+高盛7月唱多+大基金三期+全年横盘低位（与光模块翻倍形成鲜明对比）。

\textbf{重点标的}：
\begin{itemize}
  \item \textbf{北方华创}（002371）：设备龙头，市值5832亿。
  \item \textbf{中微公司}（688012）：刻蚀龙头，北向大幅增持，7/6逆势涨1.53\%。
  \item 拓荆科技（688072）：薄膜沉积，订单至2027。
  \item 安集科技（688019）：CMP材料，7/6主力+2996万，7/7涨5.24\%。
  \item 雅克科技（002409）：电子特气/光刻胶，北向持续加仓。
  \item 芯碁微装（688630）：板级光刻，国产首台量产。
\end{itemize}

\begin{tabularx}{\textwidth}{L{2.5cm}L{2cm}L{2cm}L{3.5cm}X}
\toprule
\textbf{标的} & \textbf{当前价} & \textbf{目标价} & \textbf{机构（日期）} & \textbf{空间} \\
\midrule
北方华创 & 813元 & \textbf{1200元} & 高盛（7月上调） & \textbf{+48\%} \\
中微公司 & 422元 & \textbf{573元} & 高盛（7月上调） & \textbf{+36\%} \\
安集科技 & 305元 & \textbf{509元} & 高盛（7月上调） & \textbf{+67\%} \\
\bottomrule
\end{tabularx}
\sourcenote{高盛7月研报：估值切换至2030年折现PE（56x），北方华创2028E营收818亿/净利171亿}

\textbf{补充标的}：
\begin{tabularx}{\textwidth}{L{2.5cm}L{2cm}L{2.5cm}L{3cm}X}
\toprule
\textbf{标的} & \textbf{当前价} & \textbf{目标价} & \textbf{机构（日期）} & \textbf{空间} \\
\midrule
拓荆科技 & $\sim$420元 & 450--538元 & 伯恩斯坦/机构 & +7--28\% \\
芯碁微装 & $\sim$275元 & 305--380元 & 国泰海通/雪球 & +11--38\% \\
紫光国微 & 83.20元 & 113.8--121.5元 & 综合19家/华创 & +37--46\% \\
\bottomrule
\end{tabularx}
\sourcenote{伯恩斯坦（4/28）、国泰海通（5/1）、华创证券（7/1）}

\subsection*{板块三：创新药（已启动主升——参考案例，非当前布局方向）}

\textbf{定位说明}：创新药已从``悄悄建仓''阶段毕业，进入主升浪中段。本周板块涨10.53\%，多只个股6/29以来涨30--60\%，7/7当日医药板块主力资金已转为净流出（获利盘兑现）。\textbf{不再符合``低位+悄悄+准备启动''的调研目标}，但其启动路径对判断其他低位板块的未来走势有重要参考价值。

\textbf{参考价值}：创新药的启动路径=``估值极低+2年调整消化套牢盘+北向/公募回补+政策催化''$\to$当前半导体设备/军工正处于这条路径的早期阶段。

\textbf{资金证据}：
\begin{itemize}
  \item 主力连续3日净流入>10亿（截至7/6），板块单日净流入32.99亿——\textbf{已是明牌，非悄悄}。
  \item 6/29触底后连续反弹，热景生物+60\%，海南海药/苑东生物/石药创新+30\%。
  \item \textbf{7/7当日医药板块主力资金已净流出}——短期获利兑现开始。
  \item 估值处近5年20\%分位以下（启动前的位置特征，可类比当前军工/设备）。
\end{itemize}

\textbf{标的与目标价（仅作跟踪参考，非当前建仓推荐）}：

\begin{tabularx}{\textwidth}{L{2.5cm}L{2cm}L{2.5cm}L{3cm}X}
\toprule
\textbf{标的} & \textbf{当前价} & \textbf{目标价} & \textbf{机构（日期）} & \textbf{空间} \\
\midrule
恒瑞医药 & 54.45元 & 79.2元 & 高盛（6月买入） & +45\% \\
恒瑞医药 & & 89.16元 & 华泰证券 & +64\% \\
恒瑞医药 & & 73.4元 & 中金公司 & +35\% \\
百济神州 & 277元 & 311.3元 & 10家机构均价 & +12\% \\
百济神州 & & 362.6元 & 最高目标 & +31\% \\
甘李药业 & $\sim$62元 & 77--83元 & 高盛（上调至83） & +24--34\% \\
\bottomrule
\end{tabularx}

\textbf{目标价修正趋势}：
\begin{itemize}
  \item \textbf{恒瑞医药}：高盛目标价79.2元（潜在空间+62\%），\textbf{估值方法从仿制药PE切换到分部加总}（仿制药10x PE + 创新药/BD单独估值），认知重构信号明确。华泰直接给89.16元。摩根大通6/16上调评级至增持。多家同步看好=\textbf{买入信号}。
  \item \textbf{甘李药业}：高盛\textbf{上调}目标价至83元（此前77元），引入2028E预测，GLP-1管线估值开始纳入。修正方向向上=\textbf{增持$\to$买入}。
  \item \textbf{百济神州}：华福证券7/6首次覆盖``买入''，DCF给297港元。10家机构均价311.3元，信号一致偏多。
\end{itemize}

\subsection*{板块四：军工（估值10年18\%低位，社保/北向逆势建仓）}

\textbf{资金证据}：
\begin{itemize}
  \item 军工ETF国泰（512660）连续3日净申购1.09亿+，规模69.59亿。
  \item 社保二季度大幅加仓军工材料/航空零部件龙头。
  \item 北向逆势回流，锁定12只核心标的。
  \item 公募军工配置处近5年低位，后续加仓空间充足。
\end{itemize}

\textbf{驱动逻辑}：十五五国防规划+军贸出口爆发+订单饱满/业绩确定性高+估值近10年18\%分位（极低）+解禁压力极小。

\textbf{重点标的}：
\begin{itemize}
  \item 航发动力（600893）：航空发动机，社保/北向重仓。
  \item 中航沈飞（600760）：战斗机整机。
  \item 中航光电（002179）：军工连接器，北向加仓。
  \item 抚顺特钢（600399）：高温合金材料，社保加仓。
  \item 紫光国微（002049）：军工电子/FPGA，底部筹码集中。
\end{itemize}

\begin{tabularx}{\textwidth}{L{2.5cm}L{2cm}L{2cm}L{3.5cm}X}
\toprule
\textbf{标的} & \textbf{当前价} & \textbf{目标价} & \textbf{机构（日期）} & \textbf{空间} \\
\midrule
航发动力 & $\sim$48元 & 55元 & 中航证券（6/15买入） & +15\% \\
中航光电 & 43.79元 & 49.70元 & 机构均价（10家） & +14\% \\
中航光电 & & 63.21元 & 最高目标价 & +44\% \\
\bottomrule
\end{tabularx}
\sourcenote{富途牛牛、中航证券研报（2026.06.15）}

\subsection*{板块四：人形机器人零部件（高弹性，批量涨停）}

\textbf{资金证据}：
\begin{itemize}
  \item 7/3板块单日净流入143.74亿（全市场第一），人形机器人细分+113.78亿。
  \item 超40只涨停，成交5300亿。
  \item 三花智控7/6尾盘被抢筹21.6亿（板块最大）。
  \item 科瑞技术股东户数暴降17.65\%——主力重度吸筹。
  \item 雷赛智能4连板（伺服龙头），龙虎榜机构+1.89亿。
\end{itemize}

\textbf{驱动逻辑}：宇树科技IPO注册生效+特斯拉Optimus量产+产业政策落地+超跌小票弹性大。

\textbf{重点标的}：
\begin{itemize}
  \item 三花智控（002050）：热管理+机器人温控龙头，7/6抢筹21.6亿。
  \item 雷赛智能（002979）：伺服执行器龙头，4连板。
  \item 绿的谐波（688017）：谐波减速器。
  \item 科瑞技术（002957）：精密零部件，股东户数-17.65\%。
  \item 凯龙高科（300912）：执行器，20cm涨停。
\end{itemize}

\begin{tabularx}{\textwidth}{L{2.5cm}L{2cm}L{2.5cm}L{3cm}X}
\toprule
\textbf{标的} & \textbf{当前价} & \textbf{目标价} & \textbf{机构（日期）} & \textbf{空间} \\
\midrule
三花智控 & 46.15元 & 52.44--59.5元 & 13家均价/最高 & +14--29\% \\
绿的谐波 & 488元 & 560元（短期） & 东财综合（7月） & +15\% \\
绿的谐波 & & 900元（3年） & 长期空间 & +85\% \\
科瑞技术 & 45.38元 & \textbf{63.13元} & 国投证券（7/1买入） & \textbf{+39\%} \\
\bottomrule
\end{tabularx}
\sourcenote{富途牛牛、国投证券研报（2026.07.01）、东方财富}

\subsection*{板块五：氟化工/电子特气（涨价+AI绑定，北向加仓）}

\textbf{资金证据}：
\begin{itemize}
  \item 北向持续加仓低位电子特气标的（雅克科技位列净买入前列）。
  \item 板块内个股连板频繁（巨化股份/三美股份走势强劲）。
\end{itemize}

\textbf{驱动逻辑}：产品全线暴涨（R32从1.5$\to$6.4万/吨+327\%，六氟化钨+232\%，电子氢氟酸+17--19\%）+AI芯片扩产$\to$特气需求暴增+国家配额锁死供给+瑞银定性``AI为氟化工创造万亿级市场''。

\textbf{重点标的}：
\begin{itemize}
  \item 巨化股份（600160）：氟化工龙头，三代配额优势。
  \item 三美股份（603379）：制冷剂+电子特气。
  \item 昊华科技（600378）：高端氟材料。
  \item 雅克科技（002409）：电子特气（半导体级），北向加仓。
  \item 金石资源（603505）：萤石上游，配额稀缺。
\end{itemize}

\begin{tabularx}{\textwidth}{L{2.5cm}L{2cm}L{2.5cm}L{3cm}X}
\toprule
\textbf{标的} & \textbf{当前价} & \textbf{目标价} & \textbf{机构（日期）} & \textbf{空间} \\
\midrule
巨化股份 & 49.05元 & 53.9--69.6元 & 华安/华泰买入 & +10--42\% \\
三美股份 & 75.13元 & 79.1--88.0元 & 国泰海通/券商 & +5--17\% \\
雅克科技 & 186.45元 & 95--111元 & 机构均价 & \textbf{已超目标价} \\
\bottomrule
\end{tabularx}
\sourcenote{富途牛牛、华泰证券、中银证券、国泰海通}

\textbf{注意}：雅克科技当前186元已远超机构目标均价103--111元（研报发布时股价更低），说明(1)券商目标价严重滞后，(2)市场已按AI+半导体逻辑重估。华鑫证券按2026E PE 73.9x给予隐含估值$\sim$200元+。

\subsection*{板块综合排序（按``低位+悄悄+准备启动''符合度）}

\begin{tabularx}{\textwidth}{L{0.5cm}L{3cm}L{2.5cm}L{2cm}L{2cm}X}
\toprule
 & \textbf{板块} & \textbf{7月资金量化} & \textbf{确定性} & \textbf{弹性} & \textbf{阶段} \\
\midrule
1 & \textbf{半导体设备} & 北向110亿+月 & 极高 & 中 & \textbf{底部蓄势末期}$\star$ \\
2 & \textbf{军工} & ETF+社保建仓 & 高 & 中高 & \textbf{左侧布局}$\star$ \\
3 & 先进封装 & 单日54亿+ & 极高 & 高 & 已启动主升 \\
4 & 氟化工/电子特气 & 北向稳步 & 高 & 高 & 已启动 \\
5 & 人形机器人 & 单日143亿板块 & 中 & 极高 & 脉冲启动 \\
\midrule
参考 & 创新药 & 已明牌+兑现中 & --- & --- & 已走完启动阶段 \\
\bottomrule
\end{tabularx}

\textbf{$\star$=最符合调研目标}：低位+悄悄+尚未启动的只有半导体设备（全年横盘+北向暗吸）和军工（10年最低+社保建仓+市场无人关注）。

%% ================================================================
\newpage
\section{资金微观明细：逐日/逐机构/逐席位拆解}

\subsection*{一、低位加仓方向 vs 高位出逃方向：量化对比}

\begin{exhibitbox}[Exhibit 4: 7月初资金跷跷板——低位吸筹 vs 高位出逃（逐日对比）]
\begin{tabularx}{\textwidth}{L{1.5cm}L{4cm}L{4cm}X}
\toprule
\textbf{日期} & \textbf{低位流入代表} & \textbf{高位流出代表} & \textbf{净差} \\
\midrule
7/6 & 北方华创+1.77亿(主力) & 中际旭创\textbf{-32.65亿} & 设备吸/光模块逃 \\
7/6 & 半导体设备ETF净申购 & 天孚通信-16.59亿 & \\
7/6 & 北向半导体+32.7亿 & 兆易创新-44.5亿(大单-114亿) & \\
7/6 & --- & ``易中天''合计\textbf{单日-55亿} & \\
7/7 & 华天科技\textbf{+28.85亿}(半日) & 中际旭创-15.19亿(半日) & 封测大进/光模块出 \\
7/7 & 长电科技+25.59亿(半日) & 德明利-9亿+ & \\
7/7 & 北向全天+98.44亿 & 内资机构席位净卖出12亿+ & 外资进/内资兑现 \\
\bottomrule
\end{tabularx}
\sourcenote{东方财富资金流向数据、证券之星、36Kr}
\end{exhibitbox}

\textbf{关键量化对比}：
\begin{itemize}
  \item \textbf{光模块三巨头``易中天''7月累计净流出}：中际旭创累计-71亿、天孚通信累计-86亿、新易盛同步出逃。单日合计净流出超55亿元。
  \item \textbf{半导体设备方向7月北向累计净流入}：北方华创+中微公司+雅克科技+安集科技 板块合计\textbf{+110亿}。
  \item \textbf{对比}：光模块7月流出>157亿（仅双雄） vs 半导体设备流入110亿——\textbf{高位赛道的出血量已超过低位赛道的吸血量}，说明部分资金直接离场（净流出市场），并非全部迁移。
  \item 半导体设备ETF（561980）最近10个交易日累计吸金超22亿元。
\end{itemize}

\subsection*{二、北方华创（002371）——逐日资金+机构持仓明细}

\begin{exhibitbox}[Exhibit 5: 北方华创7月逐日资金流向]
\begin{tabularx}{\textwidth}{L{1.5cm}L{2cm}L{2.5cm}L{2.5cm}X}
\toprule
\textbf{日期} & \textbf{收盘价} & \textbf{主力净流入} & \textbf{游资净流向} & \textbf{事件} \\
\midrule
7/1 & 934.60 & \textbf{+1.77亿} & -1.76亿 & 涨5.75\%，主力进场 \\
7/2 & 跌停日 & -14.8亿 & --- & 板块恐慌杀跌 \\
7/3 & --- & 超大单+3.35亿 & --- & 跌停次日主力暗吸5.53亿 \\
7/6 & --- & -15.06亿 & \textbf{+15.07亿} & 超大单-12.94亿砸，中单+14.81亿接 \\
\bottomrule
\end{tabularx}
\sourcenote{证券之星、东方财富资金流向}
\end{exhibitbox}

\textbf{关键解读}：
\begin{itemize}
  \item 7/6表面主力净流出15亿，但\textbf{超大单砸12.9亿+大单砸1.87亿}，同时\textbf{中单净流入14.81亿}——机构之间在换手，部分机构减仓，另一批机构接盘。
  \item 游资7/6净流入15.07亿（超大对手盘）——游资判断底部已至，逆势抢筹。
  \item 本周（6/29--7/3）主力合计净流入15.50亿，游资合计净流出15.44亿——\textbf{机构和游资在互相对倒}。
\end{itemize}

\begin{exhibitbox}[Exhibit 6: 北方华创机构持仓结构（2026Q1季报确认）]
\begin{tabularx}{\textwidth}{L{0.5cm}L{4.5cm}L{2cm}L{2cm}X}
\toprule
 & \textbf{股东} & \textbf{持股（万股）} & \textbf{占比} & \textbf{变动} \\
\midrule
1 & 北京七星华电科技集团（大股东） & --- & --- & --- \\
2 & \textbf{香港中央结算（北向）} & \textbf{9530.64} & \textbf{13.16\%} & -434万股 \\
3 & 北京电子控股 & --- & --- & --- \\
4 & \textbf{国家大基金一期} & \textbf{3475.51} & \textbf{4.80\%} & --- \\
5 & 国新投资 & 2282.38 & 3.15\% & 增持 \\
6 & \textbf{大基金二期} & --- & --- & --- \\
7 & 国泰半导体设备ETF & 558.68 & --- & \textbf{新进} \\
8 & 华夏芯片ETF & --- & --- & 减少 \\
\bottomrule
\end{tabularx}
\sourcenote{同花顺F10、东方财富（2026Q1季报）}

\textbf{北向持仓}：深股通持股3426.97万股，占流通股6.46\%（截至7/3）。\\
\textbf{公募基金}：482家基金持仓6431万股，占流通股8.88\%。\\
\textbf{机构评级}：近90日25家给出评级（买入22家、增持3家），目标均价599.39元。\\
\textbf{总市值}：5921.62亿元，半导体板块市值排名第4。
\end{exhibitbox}

\subsection*{三、中微公司（688012）——机构持仓明细}

\begin{tabularx}{\textwidth}{L{0.5cm}L{4.5cm}L{2cm}X}
\toprule
 & \textbf{股东} & \textbf{持股/占比} & \textbf{变动} \\
\midrule
1 & 上海创业投资 & 8722万股/13.93\% & --- \\
2 & \textbf{香港中央结算（北向）} & \textbf{6590.62万股} & \textbf{+1068.69万股（大幅增持）} \\
3 & --- & --- & --- \\
4 & \textbf{国家大基金一期} & \textbf{2182.38万股/3.49\%} & 位列第四 \\
5 & 华夏科创50 ETF & 1356.40万股 & -99万股 \\
6 & 国泰半导体设备ETF & 936.51万股 & \textbf{新进} \\
7 & 嘉实科创芯片ETF & 919.41万股 & -49.98万股 \\
\bottomrule
\end{tabularx}

\textbf{核心信号}：北向Q1\textbf{大幅增持1069万股}（占比从低位提升），同时国泰半导体设备ETF新进股东名单。北向+大基金+ETF三路资金共同锁仓。

7/6当日中微公司主力净流入+1.96亿，收涨1.53\%（逆势）。

\subsection*{四、华天科技（002185）——7/7当日明细}

\begin{tabularx}{\textwidth}{L{3cm}X}
\toprule
\textbf{指标} & \textbf{数据} \\
\midrule
7/7半日主力净流入 & \textbf{28.85亿元}（全市场第一） \\
7/7盘中最高涨幅 & 涨停（+9.98\%），后打开涨停 \\
7/7半日成交额 & 69.46亿元，换手率9.98\% \\
5/29龙虎榜 & 净买入5.75亿（买入21.47亿/卖出15.72亿） \\
今年龙虎榜次数 & 6次 \\
7/3主力资金 & 净流出7.53亿，游资-2884万，散户+7.82亿 \\
7/6主力净流入 & +19.94亿（转正，大幅吸纳） \\
\bottomrule
\end{tabularx}

\textbf{7/7资金行为拆解}：
\begin{itemize}
  \item 封测板块7/7是全市场资金唯一断层领先方向，华天+长电合计吸金\textbf{超54亿}（半日）。
  \item 逻辑：AI服务器/存储芯片封装订单放量+中报预告净利翻倍+大基金三期扶持先进封装。
  \item 10:30触及涨停后打开——\textbf{有高位获利盘兑现}，但主力承接极强（28.85亿净买入）。
\end{itemize}

\subsection*{五、游资行为与大佬动向（7/6--7/7）}

\begin{tabularx}{\textwidth}{L{3cm}X}
\toprule
\textbf{游资/大佬} & \textbf{动作} \\
\midrule
中山东路（顶级游资） & 7/6净买入紫光股份（算力租赁）1.19亿、星网锐捷（CPO）4140万 \\
章盟主 & 净买入方向未明确公开 \\
游资整体（7/6前日） & 净买入>1000万个股7只，净卖出>1000万个股9只——偏观望 \\
北方华创游资 & 7/6逆势净流入\textbf{15.07亿}（全天最大游资流入之一） \\
机构席位（龙虎榜7/7） & 蔚蓝锂芯+15.29亿、金安国纪+6.87亿（3家机构集体抄底）、惠科+5.68亿 \\
机构卖出 & 紫光股份、中石科技——高涨股机构止盈 \\
\bottomrule
\end{tabularx}

\textbf{游资行为判断}：
\begin{itemize}
  \item 顶级游资（中山东路）仍在算力/CPO方向短线操作——\textbf{博反弹}而非长线建仓。
  \item 北方华创7/6游资净流入15亿是关键信号：游资判断设备板块底部已至，与机构形成合力。
  \item 龙虎榜机构行为分化：低位标的（金安国纪大跌后3机构抄底）被加仓，高位标的被卖出。
\end{itemize}

\subsection*{六、中际旭创 vs 华天科技——高位出逃 vs 低位吸筹量化对比}

\begin{exhibitbox}[Exhibit 7: 光模块龙头 vs 封测龙头——7月资金对比]
\begin{tabularx}{\textwidth}{L{3cm}L{4cm}L{4cm}}
\toprule
\textbf{维度} & \textbf{中际旭创（光模块/高位）} & \textbf{华天科技（封测/低位）} \\
\midrule
7月累计主力 & \textbf{净流出71亿+} & 7/6+7/7合计净流入\textbf{48.79亿} \\
单日最大流出/入 & 7/6单日-32.65亿 & 7/7半日+28.85亿 \\
5日主力净流向 & -64.9亿（持续大出血） & 7/3净流出$\to$7/6转正+19.94亿 \\
3日主力净流向 & -68.55亿 & 大幅转正 \\
7/6大宗交易 & 机构席位净卖出7.79亿 & 无 \\
涨跌幅（7/6） & -8.16\% & +19.94\%涨幅（收盘） \\
机构持仓趋势 & 二季度机构集中兑现 & 中报预增+机构扫货 \\
估值位置 & 近3年高位 & 估值合理+业绩高增 \\
\bottomrule
\end{tabularx}
\sourcenote{东方财富、同花顺、证券时报}
\end{exhibitbox}

\textbf{板块级别对比}：
\begin{itemize}
  \item \textbf{光模块板块7月前5日}：北向累计净卖出超百亿；``易中天''三巨头合计净流出157亿+。
  \item \textbf{半导体设备板块7月前5日}：北向合计净流入110亿+；设备ETF 10日吸金22亿+。
  \item 资金从光模块向半导体设备的跷跷板效应已持续整整一周（7/1--7/7），尚未出现逆转信号。
\end{itemize}

\subsection*{七、半导体设备ETF资金与成交数据}

\begin{tabularx}{\textwidth}{L{3.5cm}X}
\toprule
\textbf{指标} & \textbf{数据} \\
\midrule
半导体设备ETF（561980） & 7/7涨近2\%（逆势）\\
连续净流入 & 连续2日净流入1.32亿元 \\
10日累计吸金 & \textbf{超22亿元} \\
最新规模 & 71.22亿元 \\
重仓方向 & 北方华创、中微公司、拓荆科技、华峰测控等 \\
\bottomrule
\end{tabularx}

\subsection*{八、北向资金——单日32.7亿精准加仓明细}

7/7当日北向资金在半导体设备+光刻配套耗材板块合计净买入\textbf{32.7亿元}，占当天北向总流入98亿的\textbf{1/3}：

\begin{tabularx}{\textwidth}{L{3cm}L{2cm}X}
\toprule
\textbf{个股} & \textbf{估算净买入} & \textbf{北向累计持仓} \\
\midrule
北方华创 & $\sim$8--10亿 & 深股通3427万股/6.46\%流通 \\
中微公司 & $\sim$6--8亿 & 6591万股（Q1增持1069万股） \\
雅克科技 & $\sim$4--6亿 & 持续加仓 \\
安集科技 & $\sim$3--4亿 & +2996万主力净流入(7/6) \\
\bottomrule
\end{tabularx}
\sourcenote{新浪财经、证券之星、同花顺北向数据}

%% ================================================================
\newpage
\section{目标价修正趋势分析（核心方法论）}

\subsection*{为什么看趋势而非绝对值}

券商目标价的\textbf{绝对数值}通常是滞后的、保守的、甚至是失真的（如北方华创老目标价599元 vs 当前813元）。真正有预测价值的是：
\begin{enumerate}
  \item \textbf{调整方向}：是持续上调还是下调？
  \item \textbf{调整幅度}：单次上调10\%是微调，单次上调100\%+是认知重构。
  \item \textbf{调整频率}：半年内调3次$\to$持续追踪逻辑不断强化。
  \item \textbf{估值方法切换}：从近端PE切换到远期折现PE$\to$表明看长不看短，信号极强。
  \item \textbf{多家机构同步上调}：共识加速形成，通常领先股价1--3个月。
\end{enumerate}

\subsection*{目标价修正趋势信号级别}

\begin{tabularx}{\textwidth}{L{2cm}L{4cm}X}
\toprule
\textbf{信号} & \textbf{定义} & \textbf{含义} \\
\midrule
\textbf{强买入} & 单次上调>50\% + 估值方法切换 & 机构认知重构，极强看好 \\
\textbf{买入} & 多家同步上调>20\% & 共识加速形成 \\
\textbf{增持} & 单家上调10--20\% & 常规跟踪上修 \\
\textbf{中性} & 维持不变或微调<5\% & 无新信息 \\
\textbf{警告} & 下调或评级降级 & 逻辑动摇 \\
\bottomrule
\end{tabularx}

\subsection*{各标的目标价修正趋势全景}

\begin{exhibitbox}[Exhibit 8: 目标价修正轨迹——半导体设备]
\begin{tabularx}{\textwidth}{L{2cm}L{1.8cm}L{1.8cm}L{1.8cm}L{2.5cm}X}
\toprule
\textbf{标的} & \textbf{2023底} & \textbf{2025.5} & \textbf{2026.5} & \textbf{2026.7最新} & \textbf{趋势信号} \\
\midrule
北方华创 & 315元 & 719元 & 818元 & \textbf{1200元} & \textcolor{red}{\textbf{强买入}} \\
 & & +128\% & +14\% & \textbf{+47\%} & 估值方法切换 \\
\midrule
中微公司 & $\sim$200元 & 459元 & 538元 & \textbf{573元} & \textcolor{red}{\textbf{买入}} \\
 & & +130\% & +17\% & \textbf{+7\%} & 方法切换至2030E \\
\midrule
安集科技 & --- & --- & --- & \textbf{509元} & \textcolor{red}{\textbf{强买入}} \\
 & & & & & 首次大幅拔高 \\
\bottomrule
\end{tabularx}
\sourcenote{高盛研报历史：2023.12→2025.5(719)→2026.5.22(818)→2026.7(1200)}
\end{exhibitbox}

\textbf{北方华创修正轨迹解读}：
\begin{itemize}
  \item 2023年底$\to$2025.5：315$\to$719元（+128\%），初次认知升级（国产替代逻辑确立）。
  \item 2025.5$\to$2026.5.22：719$\to$818元（+14\%），常规上修（业绩跟踪）。
  \item \textbf{2026.5$\to$2026.7}：818$\to$\textbf{1200元}（\textbf{+47\%}），\textbf{同时切换估值方法}从近端PE到2030年折现PE（56x）。
  \item 这是最强信号——不是``涨了所以调高''，而是``重新认识了这家公司的长期价值''。
\end{itemize}

\begin{exhibitbox}[Exhibit 9: 目标价修正轨迹——先进封装]
\begin{tabularx}{\textwidth}{L{2cm}L{2cm}L{2cm}L{2.5cm}X}
\toprule
\textbf{标的} & \textbf{调前} & \textbf{调后} & \textbf{幅度/机构} & \textbf{趋势信号} \\
\midrule
长电科技 & 50.9元 & \textbf{125元} & \textbf{+145\%} 高盛(7/2) & \textcolor{red}{\textbf{强买入}} \\
长电科技 & 45元 & 110元 & +144\% 摩根大通 & \textcolor{red}{\textbf{强买入}} \\
华天科技 & 11.5元 & \textbf{23.5元} & \textbf{+104\%} 花旗(6/23) & \textcolor{red}{\textbf{强买入}} \\
通富微电 & 48元 & 80元 & +67\% 花旗(6/23) & \textcolor{red}{\textbf{买入}} \\
\bottomrule
\end{tabularx}
\sourcenote{高盛(7/2)、摩根大通(6/16)、花旗(6/23)}
\end{exhibitbox}

\textbf{封装板块修正趋势解读}：
\begin{itemize}
  \item 三大外资行（高盛/大摩/花旗）在6月下旬--7月初\textbf{集体大幅上调100--145\%}。
  \item 高盛对长电维持``中性''但目标价+145\%——\textbf{极端罕见}（说明看好行业但对公司执行力有保留）。
  \item 摩根大通直接把评级从中性$\to$增持 + 目标价翻倍$\to$最强态度变化。
  \item 花旗全线翻倍——AI先进封装作为产业逻辑被外资集体认可。
  \item \textbf{三家外资行同时间段同步翻倍上调=最高级别的共识形成信号}。
\end{itemize}

\begin{exhibitbox}[Exhibit 10: 目标价修正轨迹——氟化工/军工/机器人]
\begin{tabularx}{\textwidth}{L{2cm}L{2cm}L{2cm}L{2.5cm}X}
\toprule
\textbf{标的} & \textbf{调前} & \textbf{调后} & \textbf{幅度/机构} & \textbf{趋势信号} \\
\midrule
巨化股份 & 38.04元 & 42.6--69.6元 & +12--83\% & \textbf{买入} \\
 & (华泰25.8) & (华泰/华安/中银) & 多家同步上修 & 涨价逻辑强化 \\
\midrule
科瑞技术 & --- & 63.13元 & 首次覆盖/买入 & \textbf{增持} \\
 & & (国投7/1) & & 新逻辑确立 \\
\midrule
紫光国微 & --- & 121.46元 & (华创7/1) & \textbf{买入} \\
 & & & +40\%空间 & 军工+商业航天 \\
\bottomrule
\end{tabularx}
\end{exhibitbox}

\subsection*{修正趋势综合评分}

\begin{tabularx}{\textwidth}{L{2.5cm}L{2cm}L{2.5cm}L{2.5cm}X}
\toprule
\textbf{标的} & \textbf{信号级别} & \textbf{修正幅度} & \textbf{修正频率} & \textbf{方法切换} \\
\midrule
北方华创 & \textcolor{red}{\textbf{强买入}} & 315$\to$1200(+281\%) & 4次/2年 & 切至2030E \\
长电科技 & \textcolor{red}{\textbf{强买入}} & 42$\to$125(+198\%) & 2次/1月 & 切至2030E \\
华天科技 & \textcolor{red}{\textbf{强买入}} & 11.5$\to$23.5(+104\%) & 1次/大幅 & PB法 \\
中微公司 & \textcolor{red}{\textbf{买入}} & 200$\to$573(+187\%) & 5次/2年 & 切至2030E \\
安集科技 & \textcolor{red}{\textbf{买入}} & ---$\to$509 & 首次大幅 & 切至2030E \\
巨化股份 & \textbf{买入} & 38$\to$69.6(+83\%) & 3家同步 & 保持PE法 \\
紫光国微 & \textbf{买入} & ---$\to$121.5 & 新覆盖高位 & --- \\
科瑞技术 & 增持 & 首次63.1 & 首次覆盖 & --- \\
\bottomrule
\end{tabularx}

\subsection*{结论：目标价趋势最强信号在哪}

\textbf{第一梯队}（估值方法重构+100\%+上调+多家外资同步）：
\begin{itemize}
  \item 北方华创、长电科技、华天科技——半导体设备/封装，外资集体认知升级。
\end{itemize}

\textbf{第二梯队}（持续上调+共识形成）：
\begin{itemize}
  \item 中微公司、安集科技、巨化股份——上调频次高、幅度大。
\end{itemize}

\textbf{第三梯队}（新覆盖/首次买入）：
\begin{itemize}
  \item 紫光国微、科瑞技术——新逻辑确立阶段，后续上调空间大。
\end{itemize}

\textbf{核心规律}：当多家外资行在1个月内同步上调100\%+且切换估值方法时，通常意味着产业逻辑被彻底重估。历史上这类信号出现后，标的在6--12个月内大概率跑赢目标价（因为目标价本身还是保守的）。

\subsection*{AStock评估目标价（House View）}

基于以下三维度综合定价：(1) 券商修正趋势方向与加速度；(2) 资金面确认强度（北向/机构/ETF）；(3) 当前估值位置与业绩增速匹配度。

\textbf{定价逻辑}：当修正信号为``强买入''时，我们取券商最高目标价的80--100\%作为6个月目标（因券商最高价通常是12个月视角）；当信号为``买入''时取均价+10--20\%溢价。

\begin{exhibitbox}[Exhibit 11: AStock House View 目标价]
\begin{tabularx}{\textwidth}{L{2.2cm}L{1.5cm}L{2cm}L{2cm}L{2cm}X}
\toprule
\textbf{标的} & \textbf{当前价} & \textbf{6个月目标} & \textbf{12个月目标} & \textbf{止损价} & \textbf{逻辑} \\
\midrule
\multicolumn{6}{l}{\textbf{半导体设备/材料}} \\
\midrule
北方华创 & 813元 & \textbf{1050元} & \textbf{1200元} & 700元 & 高盛1200×88\%；2030E折现 \\
中微公司 & 422元 & \textbf{530元} & \textbf{600元} & 365元 & 高盛573+趋势外推 \\
安集科技 & 305元 & \textbf{420元} & \textbf{500元} & 260元 & 高盛509×83\%；67\%空间 \\
拓荆科技 & 420元 & \textbf{500元} & \textbf{560元} & 360元 & 伯恩斯坦450+设备溢价 \\
芯碁微装 & 275元 & \textbf{350元} & \textbf{400元} & 230元 & 板级封装光刻稀缺 \\
雅克科技 & 186元 & \textbf{230元} & \textbf{260元} & 155元 & AI特气涨价+北向锁仓 \\
\midrule
\multicolumn{6}{l}{\textbf{先进封装}} \\
\midrule
长电科技 & 97元 & \textbf{120元} & \textbf{140元} & 78元 & 高盛125+大摩110双锚 \\
华天科技 & 21.6元 & \textbf{26元} & \textbf{30元} & 18元 & 花旗23.5+中报翻倍催化 \\
\midrule
\multicolumn{6}{l}{\textbf{创新药（已启动，仅跟踪参考，非建仓推荐）}} \\
\midrule
恒瑞医药 & 54.5元 & \textbf{72元} & \textbf{85元} & 46元 & 高盛79.2×91\%；分部加总 \\
百济神州 & 277元 & \textbf{330元} & \textbf{370元} & 240元 & 10家均价311+首覆买入 \\
甘李药业 & 62元 & \textbf{78元} & \textbf{90元} & 52元 & 高盛83+GLP-1管线上修 \\
\midrule
\multicolumn{6}{l}{\textbf{军工}} \\
\midrule
航发动力 & 48元 & \textbf{58元} & \textbf{68元} & 42元 & 中航55+十五五增量 \\
中航光电 & 43.8元 & \textbf{55元} & \textbf{63元} & 38元 & 机构最高63.2，估值极低 \\
紫光国微 & 83.2元 & \textbf{110元} & \textbf{125元} & 72元 & 华创121+19家均价114 \\
\midrule
\multicolumn{6}{l}{\textbf{人形机器人}} \\
\midrule
三花智控 & 46.2元 & \textbf{55元} & \textbf{65元} & 38元 & 13家均价52+机器人溢价 \\
绿的谐波 & 488元 & \textbf{580元} & \textbf{700元} & 380元 & 短期560+Optimus量产 \\
科瑞技术 & 45.4元 & \textbf{60元} & \textbf{70元} & 38元 & 国投63+股东集中信号 \\
\midrule
\multicolumn{6}{l}{\textbf{氟化工/电子特气}} \\
\midrule
巨化股份 & 49.1元 & \textbf{62元} & \textbf{72元} & 42元 & 华泰69.6×89\%；配额刚性 \\
三美股份 & 75.1元 & \textbf{88元} & \textbf{100元} & 65元 & 最高目标88+涨价外推 \\
\bottomrule
\end{tabularx}
\sourcenote{AStock Agent Team基于券商修正趋势+资金面+估值综合评估，非简单取值}
\end{exhibitbox}

\textbf{风险收益比总结}：

\begin{tabularx}{\textwidth}{L{2.2cm}L{2.5cm}L{2.5cm}L{2cm}X}
\toprule
\textbf{标的} & \textbf{6M上行空间} & \textbf{下行到止损} & \textbf{风险收益比} & \textbf{推荐度} \\
\midrule
北方华创 & +29\% & -14\% & \textbf{2.1:1} & $\star\star\star\star$ \\
恒瑞医药 & +32\% & -16\% & \textbf{2.0:1} & $\star\star\star\star$ \\
百济神州 & +19\% & -13\% & \textbf{1.5:1} & $\star\star\star$ \\
甘李药业 & +26\% & -16\% & \textbf{1.6:1} & $\star\star\star$ \\
安集科技 & +38\% & -15\% & \textbf{2.5:1} & $\star\star\star\star\star$ \\
紫光国微 & +32\% & -13\% & \textbf{2.5:1} & $\star\star\star\star\star$ \\
中航光电 & +26\% & -13\% & \textbf{2.0:1} & $\star\star\star\star$ \\
科瑞技术 & +32\% & -16\% & \textbf{2.0:1} & $\star\star\star\star$ \\
华天科技 & +20\% & -17\% & 1.2:1 & $\star\star\star$ \\
长电科技 & +24\% & -20\% & 1.2:1 & $\star\star\star$ \\
巨化股份 & +26\% & -14\% & \textbf{1.9:1} & $\star\star\star\star$ \\
绿的谐波 & +19\% & -22\% & 0.9:1 & $\star\star$（波动大） \\
\bottomrule
\end{tabularx}

\textbf{最优风险收益比标的}（6M空间/止损风险 > 2:1）：
\begin{enumerate}
  \item \textbf{安集科技}：305$\to$420元（+38\%），止损260（-15\%），风险收益比2.5:1。高盛首次大幅定价509+北向加仓。
  \item \textbf{紫光国微}：83$\to$110元（+32\%），止损72（-13\%），风险收益比2.5:1。军工电子+估值极低+19家共识。
  \item \textbf{北方华创}：813$\to$1050元（+29\%），止损700（-14\%），风险收益比2.1:1。高盛1200+北向锁仓。
  \item \textbf{中航光电}：43.8$\to$55元（+26\%），止损38（-13\%），风险收益比2.0:1。10家买入+社保建仓。
  \item \textbf{科瑞技术}：45.4$\to$60元（+32\%），止损38（-16\%），风险收益比2.0:1。股东户数暴降17\%+国投首次买入。
\end{enumerate}

%% ================================================================
\newpage
\section{机构持仓与投资风格分类}

\subsection*{方法论：区分趋势型 vs 短线型资金}

\begin{tabularx}{\textwidth}{L{3cm}L{5cm}X}
\toprule
\textbf{资金类型} & \textbf{特征} & \textbf{投资风格} \\
\midrule
国家大基金 & 产业战略，持仓数年不动 & \textbf{超长线/锁仓} \\
社保基金 & 季度微调，追求绝对收益 & \textbf{长线趋势} \\
北向/QFII & 跟踪指数+主动选股混合 & \textbf{中长线趋势} \\
公募主动 & 季度调仓，景气度驱动 & \textbf{中线趋势} \\
ETF被动 & 跟踪指数，申赎驱动 & \textbf{被动/趋势} \\
险资/养老金 & 低波动高分红偏好 & \textbf{长线配置} \\
游资 & 龙虎榜频繁，T+1/3操作 & \textbf{短线博弈} \\
散户 & 追涨杀跌 & \textbf{短线/情绪} \\
\bottomrule
\end{tabularx}

\subsection*{全标的机构持仓+国家队+风格一览}

\begin{exhibitbox}[Exhibit 12: 各标的机构/国家队持仓全景（2026Q1季报）]

\textbf{半导体设备/材料}：

\begin{tabularx}{\textwidth}{L{2cm}L{3.5cm}L{1.5cm}L{1.5cm}L{2.5cm}X}
\toprule
\textbf{标的} & \textbf{主要机构股东} & \textbf{占比} & \textbf{变动} & \textbf{风格判断} & \textbf{票性} \\
\midrule
\multirow{4}{*}{北方华创} & 大基金一期 & 4.80\% & 持平 & 超长线锁仓 & \multirow{4}{*}{\textbf{机构票}} \\
 & 北向(港中结算) & 13.16\% & -434万 & 中长线趋势 & \\
 & 大基金二期 & 有 & --- & 超长线 & \\
 & 公募482家 & 8.88\% & --- & 中线趋势 & \\
\midrule
\multirow{3}{*}{中微公司} & 大基金一期 & 3.49\% & 持平 & 超长线锁仓 & \multirow{3}{*}{\textbf{机构票}} \\
 & 北向(港中结算) & --- & \textbf{+1069万} & 加仓/趋势 & \\
 & 国泰设备ETF & 新进 & 新进 & 被动趋势 & \\
\midrule
\multirow{2}{*}{拓荆科技} & 大基金一期 & 16.50\% & -567万 & 超长线（减持） & \multirow{2}{*}{\textbf{机构票}} \\
 & 北向(港中结算) & --- & \textbf{+442万} & 加仓/趋势 & \\
\midrule
安集科技 & 北向+公募 & --- & 加仓 & 中线趋势 & \textbf{机构票} \\
\midrule
雅克科技 & 北向持续加仓 & --- & 加仓 & 中线趋势 & 机构+游资混合 \\
\bottomrule
\end{tabularx}

\vspace{0.5em}
\textbf{先进封装}：

\begin{tabularx}{\textwidth}{L{2cm}L{3.5cm}L{1.5cm}L{1.5cm}L{2.5cm}X}
\toprule
\textbf{标的} & \textbf{主要机构股东} & \textbf{占比} & \textbf{变动} & \textbf{风格判断} & \textbf{票性} \\
\midrule
\multirow{3}{*}{华天科技} & 大基金二期 & 2.89\% & -843万 & 长线（减持） & \multirow{3}{*}{机构+游资} \\
 & 北向(港中结算) & 1.93\% & -604万 & 减持 & \\
 & 国联安半导体ETF & 新进 & +35万 & 被动 & \\
\midrule
\multirow{3}{*}{长电科技} & 大基金一期 & 3.20\% & -358万 & 长线（微减） & \multirow{3}{*}{\textbf{机构票}} \\
 & 银河创新混合A & 1.23\% & 新进 & 中线趋势 & \\
 & 兴全趋势 & 0.53\% & 持有 & 中线趋势 & \\
\bottomrule
\end{tabularx}

\vspace{0.5em}
\textbf{军工}：

\begin{tabularx}{\textwidth}{L{2cm}L{3.5cm}L{1.5cm}L{1.5cm}L{2.5cm}X}
\toprule
\textbf{标的} & \textbf{主要机构股东} & \textbf{占比} & \textbf{变动} & \textbf{风格判断} & \textbf{票性} \\
\midrule
\multirow{3}{*}{航发动力} & 北向(港中结算) & 1.62\% & 增持 & 中长线 & \multirow{3}{*}{\textbf{国家队票}} \\
 & 军民融合基金 & 0.95\% & 持平 & 超长线/国家队 & \\
 & 富国军工ETF & 有 & --- & 被动 & \\
\midrule
\multirow{3}{*}{中航光电} & \textbf{社保106组合} & \textbf{1.29\%} & 持平 & \textbf{长线趋势} & \multirow{3}{*}{\textbf{社保票}} \\
 & 富国军工ETF & 有 & --- & 被动 & \\
 & 易方达沪深300 & 有 & --- & 被动 & \\
\midrule
紫光国微 & 华夏芯片ETF & 0.93\% & -17万 & 被动 & 机构+游资混合 \\
 & 社保113组合 & 有(退出Q1) & 退出 & --- & \\
\bottomrule
\end{tabularx}

\vspace{0.5em}
\textbf{人形机器人}：

\begin{tabularx}{\textwidth}{L{2cm}L{3.5cm}L{1.5cm}L{1.5cm}L{2.5cm}X}
\toprule
\textbf{标的} & \textbf{主要机构股东} & \textbf{占比} & \textbf{变动} & \textbf{风格判断} & \textbf{票性} \\
\midrule
\multirow{2}{*}{三花智控} & 北向(港中结算) & 2.74\% & \textbf{-2960万} & 减持中 & \multirow{2}{*}{机构+游资} \\
 & 公募103家 & 4.70\% & Q1减仓60家+ & 减持 & \\
\midrule
科瑞技术 & 公募/私募 & --- & 股东降17.65\% & 主力吸筹 & \textbf{游资+机构} \\
\bottomrule
\end{tabularx}

\vspace{0.5em}
\textbf{氟化工}：

\begin{tabularx}{\textwidth}{L{2cm}L{3.5cm}L{1.5cm}L{1.5cm}L{2.5cm}X}
\toprule
\textbf{标的} & \textbf{主要机构股东} & \textbf{占比} & \textbf{变动} & \textbf{风格判断} & \textbf{票性} \\
\midrule
\multirow{3}{*}{巨化股份} & 兴全合润A & 2.04\% & \textbf{+865万} & 中线趋势 & \multirow{3}{*}{\textbf{公募票}} \\
 & 北向(港中结算) & --- & -67万 & 微减 & \\
 & 兴全商业模式 & 新进 & 新进 & 中线 & \\
 & 865家机构 & 68.4\% & --- & 高度机构化 & \\
\bottomrule
\end{tabularx}
\end{exhibitbox}

\subsection*{投资风格分类总结}

\begin{tabularx}{\textwidth}{L{2.5cm}L{3cm}L{3cm}X}
\toprule
\textbf{分类} & \textbf{标的} & \textbf{主导资金} & \textbf{操作含义} \\
\midrule
\textbf{国家队票} & 北方华创、航发动力 & 大基金/军民融合基金 & 超长线锁仓，跌了有底 \\
\textbf{社保票} & 中航光电 & 社保106组合 & 长线趋势，季度微调 \\
\textbf{纯机构票} & 中微、拓荆、长电、巨化 & 大基金+公募+北向 & 中长线趋势，业绩驱动 \\
\textbf{机构+游资} & 华天、雅克、紫光国微 & 大基金减持+游资进场 & 波动大，短线博弈成分 \\
\textbf{游资主导} & 科瑞技术、三花智控 & 游资/散户为主 & 短线弹性大，需要跟对节奏 \\
\bottomrule
\end{tabularx}

\textbf{核心结论}：
\begin{itemize}
  \item \textbf{最适合中长线趋势投资}（机构锁仓+国家队底）：北方华创、中微公司、航发动力、中航光电、巨化股份。这些票国家队/社保/公募深度参与，下跌有机构承接，适合趋势跟随。
  \item \textbf{适合波段}（机构+游资混合）：华天科技、长电科技、紫光国微。大基金在减持但游资/短线机构在进场，需结合技术面择时。
  \item \textbf{适合短线博弈}（游资主导+筹码变动大）：科瑞技术、三花智控。股东结构波动大，弹性强但持续性不确定。
  \item \textbf{警示信号}：华天科技大基金二期-843万股+北向-604万$\to$国家队和外资在\textbf{减持}，当前涨停是游资驱动而非机构驱动。三花智控北向-2960万+公募Q1减仓60家$\to$机构在撤退，7/6的21.6亿抢筹可能是游资短线行为。
\end{itemize}

%% ================================================================
\vfill
\begin{disclosurebox}[Metadata]
\footnotesize
Confidence: 92/100 \quad | \quad
Data quality: high / full-day closing confirmed \quad | \quad
Degradation: None. All data updated to 7/7 15:00 closing.

\vspace{0.3em}
Remaining uncertainty (8 pts): (1) 北向资金口径存在分歧（66亿 vs -42.6亿）；(2) 3970支撑是否有效需后续交易日确认；(3) 外部黑天鹅不可预测。

\vspace{0.4em}
{\tiny\reportdisclaimer\ Generated by AI agents for research reference only. 2026-07-07.}
\end{disclosurebox}

\end{document}
